% HD-F1-T12. GPT-6 Astra Pro, /root/transport; 2026-11-22 UTC.
\section{HD-F1-T12: closure of the terminal rank-four tensor source}
\label{sec:f1-terminal-tensor-closure}

The terminal tensor alternative in the remaining natural-curve
source register contributes no still-missing sixfold determinant
component. More precisely, after the already proved reductions
T08--T11, every remaining period-rank-nine source maps by the
actual induction and polarized isogenies below to the scalar
augmentation source of HD-S33, and therefore has ordinary
Hermitian determinant \([-1]\). The other possible block size
has period rank at most eight by the exact source theorem
HD-S43. This closes that terminal alternative; it makes no
assertion about the separate reflection-free alternative of
moving rank at least five.

The input is the prescribed primitive factor of the original
connected full-\(G\) sphere cover \(C\to\mathbf P^1\), with
\[
 G=\mathfrak S_s\times\Gamma,\quad
 V=W\otimes_K M,\quad W=\operatorname{Aug}(K^s),\quad
 r=s-1\ge20,\quad \dim_KM=4.
 \tag{F1T12.1}
\]
Its parabolic multiplicity dimension is six, its signature is
\((3,3)\), and its period rank is nine. T11 gives
\(K=\mathbf Q(i)\) and an actual order-four reflection in
\(\Gamma\). All statements below concern this original
rank-six multiplicity space, rather than the containing
rank-\(6+d_0\) quotient used to prove T11.

\paragraph{1. An actual Gaussian block system of size one or two.}
Let \(R\lhd\Gamma\) be the subgroup generated by every
reflection of \(\Gamma\). Use an averaged positive
\(K\)-Hermitian form \(h_M\). The \(K\)-root lines of
reflections have nonorthogonality components whose spans
are mutually orthogonal. Their sum is a nonzero
\(\Gamma\)-invariant subspace and hence all of \(M\).
An invariant subspace for a component reflection group,
together with its orthogonal complement, would put each
rank-one reflection image wholly in one summand; connectedness
then puts the entire component there. Thus each span is
absolutely irreducible. Normality and irreducibility make
\(\Gamma\) permute the spans transitively. Their common
dimension divides four.

Dimensions one and two already give the required block
system. In dimension four the group \(R\) is itself an
irreducible Gaussian reflection group. The exact classification
and \(K\)-intertwiners in HD-G25 and HD-S35 give the complete
rank-four list
\[
 A_4,\ B_4,\ D_4,\ F_4,\quad
 G(4,1,4),G(4,2,4),G(4,4,4),\quad G_{29},G_{31}.
 \tag{F1T12.2}
\]
All real entries have determinant in \(\{\pm1\}\).
The same holds for \(G(4,2,4)\) and \(G(4,4,4)\), directly
from their diagonal-product condition and permutation sign.
The literal generators of \(G_{29},G_{31}\) in HD-S32 are
\(I-2vv^*\) with \(v^*v=1\), so each has determinant
\(-1\) and their whole determinant images lie in
\(\{\pm1\}\). None can contain an order-four reflection,
whose determinant is \(i\) or \(-i\).
Every reflection of \(\Gamma\) belongs to \(R\), so T11
leaves exactly \(R=G(4,1,4)\).

In this standard model the root lines of order-four
reflections are precisely the four coordinate lines.
Indeed a monomial reflection with nontrivial permutation
must interchange two coordinates, with reciprocal phases,
and has order two. Every order-four reflection is therefore
a single coordinate multiplication by \(i\) or \(-i\).
Conjugation by \(\Gamma\) preserves the set of these actual
root lines, so \(\Gamma\) preserves this coordinate line
system. The \(K\)-linear intertwiner from G25 pulls it back
to four actual \(K\)-lines in the original \(M\).
The invariant form is a positive rational scalar times the
standard form by irreducibility of \(R\), so these lines
are orthogonal for the original form as well.

We have therefore obtained, in every case, an actual
transitive orthogonal block system
\[
 M=\bigoplus_{j=1}^{t}U_j,\qquad
 a=\dim_K U_j\in\{1,2\},\qquad t=4/a .
 \tag{F1T12.3}
\]
No classification of the block-stabilizer group is needed.

\paragraph{2. Induce the original tensor, with its original field and form.}
Fix \(U=U_1\), let \(H_\Gamma\le\Gamma\) stabilize it,
and let \(\sigma:H_\Gamma\to GL_K(U)\) be the actual
restriction. It is absolutely irreducible: a proper
invariant subspace would, by its translates in the other
blocks, give a proper \(\Gamma\)-invariant subspace of \(M\).
For representatives \(t_j\) of \(\Gamma/H_\Gamma\), the
literal map is
\[
 \bigoplus_{j=1}^{t}(W\otimes U)\longrightarrow V,\qquad
       (v_j)\longmapsto
            \sum_j(1_W\otimes\rho_M(t_j))v_j .
 \tag{F1T12.4}
\]
It realizes
\[
 V=\operatorname{Ind}_{H}^{G}(W\otimes U),\qquad
 H=\mathfrak S_s\times H_\Gamma,\qquad
 \psi(\pi,h)=P_\pi|_W\otimes\sigma(h).
 \tag{F1T12.5}
\]
For example, \(\gamma t_j=t_k h\) gives precisely the
block matrix \(P_\pi\otimes\sigma(h)\). The inverse is
the direct sum of the orthogonal block projections followed
by \(1_W\otimes\rho_M(t_j)^{-1}\).

The induced trace formula and the actual \(K\)-matrices give
\[
 K=\mathbf Q(\operatorname{tr}_M\Gamma)
       \subseteq\mathbf Q(\operatorname{tr}\sigma(H_\Gamma))
       \subseteq K .
 \tag{F1T12.6}
\]
Thus the faithful image \(\Gamma_a=\sigma(H_\Gamma)\)
has exact character field \(K\). The representation \(\psi\)
has this same exact field, since its traces at \((1,h)\)
are \(r\operatorname{tr}\sigma(h)\).
It is absolutely irreducible as the external tensor of
the two absolutely irreducible representations.

The positive coefficient form on \(V\) is exactly
\(h_W\otimes h_M\) after defining
\[
 h_M(u,u')=
 \frac{h_V(w\otimes u,w\otimes u')}{h_W(w,w)}
 \quad(0\ne w\in W).
 \tag{F1T12.7}
\]
Schur averaging over the rational irreducible
\(\mathfrak S_s\)-module \(W\) proves this identity on all
pairs. The map (F1T12.4) therefore identifies the form
with the orthogonal sum of the restrictions, without an
extra scalar. This is the coefficient pairing used below.

\paragraph{3. The intermediate curves and the exact primitive corner.}
Put \(N=\ker\sigma\), viewed as \(1\times N\le H\), and
take the actual curves and finite maps
\[
 B=C/H,\qquad C_a=C/(1\times N),\qquad
 \alpha:C\longrightarrow C_a .
 \tag{F1T12.8}
\]
The kernel of \(\psi\) is exactly \(1\times N\).
An identity tensor requires both factors to be scalar,
and a standard permutation on the augmentation module
for \(s\ge21\) is scalar only for the identity permutation.
Thus
\[
 C_a\longrightarrow B\text{ is a connected full cover
 with group }\mathfrak S_s\times\Gamma_a,
 \qquad \deg\alpha=|N|,\quad [G:H]=t.
 \tag{F1T12.9}
\]
The representation on \(W\otimes U\), of rank \(ar\),
is faithful for that group.

Here is the primitive-corner identification, including its
normalization. Choose \(0\ne v\in W\otimes U\), set
\(q_v=h_V(v,v)\), and form the actual rational
matrix-coefficient idempotents
\[
 \begin{split}
 e_G&=\frac{2(4r)}{|G|q_v}
    \sum_{g\in G}\operatorname{Re}h_V(v,\rho(g^{-1})v)\,g,\\
 e_H&=\frac{2(ar)}{|H|q_v}
    \sum_{h\in H}\operatorname{Re}h_V(v,\psi(h^{-1})v)\,h .
 \end{split}
 \tag{F1T12.10}
\]
For \(g\notin H\), the two vectors belong to orthogonal
blocks and the coefficient is zero. Moreover
\((4r)/|G|=(ar)/|H|\). Hence \(e_G=e_H\) literally in
\(\mathbf Q[G]\). The latter is
\(\operatorname{av}_N\) times the inflation of the
primitive corner \(e_a\) of \(\mathfrak S_s\times\Gamma_a\).
Integer multiples clearing denominators give the actual
abelian subvarieties in both Jacobians, and
\[
 |N|^{-1}\alpha^* e_a\alpha_*=e_G,\qquad
 \alpha_*\alpha^*=[|N|],\qquad
 (\alpha^*)^*\Theta_C=|N|\Theta_{C_a}.
 \tag{F1T12.11}
\]
Thus pullback is an actual isogeny from the primitive
factor in \(JC_a\) onto the specified block corner in \(JC\).
Its determinant multiplier in \(K\)-rank six is
\(|N|^6=N_{K/\mathbf Q}(|N|^3)\).

An originally prescribed coefficient line need not lie in
the selected block. For \(v_0\) on that original line use
the actual matrix unit
\[
 T(x)=v\,\frac{h_V(v_0,x)}{h_V(v_0,v_0)},\qquad
 b=\frac{2(4r)}{|G|}
   \sum_{g\in G}\operatorname{Re}
      \operatorname{tr}(T\rho(g^{-1}))\,g .
 \tag{F1T12.12}
\]
Fourier orthogonality gives
\(b^\dagger b=(q_v/h_V(v_0,v_0))e_{v_0}\).
An integer-cleared \(Db\) is an actual \(K\)-linear
isogeny between the two primitive corners, with positive
rational polarization multiplier
\(D^2q_v/h_V(v_0,v_0)\). Its sixth power is again a norm.
These are the explicit same-corner maps of G27 and S41,
now applied to the original tensor (F1T12.1).

On local systems (F1T12.4) is the finite-sheet cellular
induction map. At a ramification cycle its fixed stalk is
the fixed space of the product of its block monodromies
on both sides, so it identifies middle extensions, not
only the open local systems. It identifies the parabolic
pairings and the rank-six Hodge structures. Equations
(F1T12.11)--(F1T12.12) identify the same period maps after
finite marking of the cover data. In particular the
primitive factor on \(C_a/B\) still has parabolic
dimension six and period rank nine.

\paragraph{4. The actual intermediate base is a sphere.}
Let \(g=g(B)\), and let \(c'_j\ge1\) be the moving ranks
of the nontrivial inertias of the faithful coefficient
system on \(C_a/B\). Irreducibility and its dual remove
both global invariant terms, giving
\[
                       6=ar(2g-2)+\sum_j c'_j .
 \tag{F1T12.13}
\]
Since \(ar\ge20\), genus at least two is impossible.
In genus one, \(\sum c'_j=6\), so the number of true
coefficient branch values is at most six. The moduli
space of that pointed elliptic base has dimension at most
six; the unpointed case is incompatible with the displayed
sum. Finite monodromy gives only finitely many choices of
the cover over a fixed pointed base. The original period
map factors through this actual cover moduli by the maps
just proved, so genus one has period rank at most six.
Therefore \(B=\mathbf P^1\). Its period rank nine requires
at least twelve nonidentity coefficient branch values.
Any ramification lost upon quotienting by \(N\) has already
been removed by the actual map \(\alpha\); its extra labels
do not change this quotient factor.

\paragraph{5. Check the full hypotheses of the two inherited source results.}
For \(a=2\), the group in (F1T12.9) is literally
\(\mathfrak S_s\times\Gamma_2\), not an unspecified
overgroup or a correlated projection. Its specified
representation is
\[
 \operatorname{Aug}(K^s)\otimes_K U,\quad
 \dim_KU=2,\quad r=s-1\ge20\ge12.
 \tag{F1T12.14}
\]
It is faithful and absolutely irreducible, has exact
character field \(K\), and uses the product positive form
already identified in (F1T12.7). It contains the entire
supported reflection subgroup
\(\mathfrak S_s\times1\), acting as the standard
rank-\(r\) reflection representation on the first factor.
Its base is the sphere, it has at least twelve true
branch values, and its parabolic dimension is exactly six.
The resulting inertia codimension sum is \(4r+6\).
Every nonidentity augmentation tensor has moving rank at
least two, exactly the bound used in S43. For a nonscalar
\(\gamma\in\Gamma_2\), its two distinct eigenvalues match
disjoint permutation eigenspaces, giving codimension at
least \(r\). For a nonidentity scalar the cycle count gives
codimension at least \(r-1\). For \(\gamma=1\), the
codimension is twice the nonzero permutation index.
No stronger moving-rank assumption is imported from the old source.
These are all hypotheses of HD-S43, paragraphs 2--5.
That complete proof supplies its own actual rank-two
quotient, complementary projector, and period bound at
most eight. Its primitive corner is the one already
identified by (F1T12.10)--(F1T12.12), so it contradicts
the original period rank nine.

For \(a=1\), the image \(\Gamma_1\le K^\times\) is finite
and has exact character field \(K=\mathbf Q(i)\), by
(F1T12.6). Hence it is exactly \(\mu_4\).
The actual cover in (F1T12.9) has full group
\(\mathfrak S_s\times\mu_4\) and representation
\((\pi,\zeta)\mapsto\zeta P_\pi|_{\operatorname{Aug}}\).
For a chosen \(K\)-basis \(u\) of the line \(U\), put
\(q=h_U(u,u)\in\mathbf Q_{>0}\). The coefficient map
\(w\mapsto w\otimes u\) carries the standard augmentation
form to \(q h_W\); this changes the parabolic determinant
by the explicit norm \(q^6=N(q^3)\).
The class-sum and matrix-coefficient \(K\)-actions are carried
by \(e_G=e_H\) and \(\alpha^*\); the scalar character
\(\sigma\) gives exactly the inverse-deck convention in S33.
Thus its base, original primitive corner, signature and
restricted-theta polarization satisfy HD-S33 with this
explicit harmless norm multiplier.
That theorem retains the coefficient discriminant \(s\)
and proves
\[
                  [\det H]=[-s^{\,r_0-2}],
 \tag{F1T12.15}
\]
where \(r_0\) is the number of nontrivial scalar
branch values in this actual quotient cover.
Odd \(r_0\) forces at most ten true branch values and
period rank at most seven; consequently the present
rank-nine family has even \(r_0\). Formula (F1T12.15)
therefore gives \([\det H]=[-1]\), since the removed
power of \(s\) is explicitly the norm of
\(s^{(r_0-2)/2}\). The two polarization multipliers in
(F1T12.11)--(F1T12.12) preserve this ordinary determinant
class for the original prescribed factor. This is the
already supplied split Gaussian component.

\paragraph{Framework and source accounting.}
The proof uses the actual quotient maps, their graph
correspondences, and the integer-cleared group-algebra
projectors. Their \(K_0\) classes and GRR realizations
are induced by these same correspondences. The
Poincare-biextension identities are
\(\widehat{\alpha_*}\lambda_{C_a}=\lambda_C\alpha^*\)
and \((\alpha^*)^*\Theta_C=|N|\Theta_{C_a}\); the
Rosati matrix-unit identity is (F1T12.12).
Their relative cohomological maps are flat and carry the
actual Hodge and variational cycle classes. Adams
operations multiply each codimension-\(j\) character
term by \(k^j\), with the GRR Todd terms retained for
any pushforward. These constructions supply the same
factor throughout.

The primary classification is Cohen (1976), Theorem (2.4),
pp.386--387; (4.4)(i), p.399; table (4.16), p.412;
Theorem (5.12), p.435. The concrete rank-four Gaussian
models and their intertwiners are the complete inherited
proofs G25, S32 and S35. The complete S33 and S43 proofs
were read again for this application, including their
local cycle maps, integral pairings, complementary
projectors and exact polarization scalars. No new
classification or new supplier family is claimed:
the new conclusion is the closure of the previously
remaining tensor alternative by these actual maps.
